\section{Replication figures: ours compared with the authors}\label{app:comparisons}
Every figure in this gallery places our replication results alongside
measurements reconstructed from the authors' released data. There are
no standalone galleries. Each caption identifies the shared settings
and any differences in training horizon, trial count, or time coordinates.
The Figure~1 comparison also appears in the main text. Released-data
reconstructions are distinct from the original published artwork and
are not additional experiments by our group~\cite{nakkiran2021deepdouble}.

\resultplot{comparison}{figure_01}{CIFAR-10 ResNet18 width sweeps (paper Figure~1). Both authors and ours use 10\%-noise CIFAR-10 ResNet18, the same ten widths and 400 displayed epochs, with one run per width. Authors' epochs follow the notebook convention; ours are recorded epochs. Test error uses the same expected noisy-test transformation.}{app:matched-figure_01}
\resultplot{comparison}{figure_02}{Train and test heatmaps (paper Figure~2): authors' and our CIFAR-10 ResNet18 train/test heatmaps at 10\% noise, on the same ten widths and displayed horizon. Authors' epoch coordinates follow the released notebook convention; ours are recorded.}{app:matched-figure_02}
\resultplot{comparison}{figure_04}{ResNet18 error versus width on CIFAR-10 and CIFAR-100, reproducing paper Figure~4. Both sources use the same ten widths and label-noise rates of 0, 10, and 20\%. The left column shows the authors at displayed epoch 400 (measurement index 399); the right column shows ours at recorded epoch 400.}{app:matched-figure_04}
\resultplot{comparison}{figure_05}{Effect of data augmentation on CIFAR-10 CNN error, reproducing paper Figure~5. Both sources use the same eleven widths and label-noise rates of 0, 10, and 20\%. The left column shows the authors at measurement index 194; the right column shows ours after 50,000 SGD updates. The author point corresponds to approximately 49,920 updates only under the unverified assumption of 256 updates per measurement, so training duration is not confirmed to match.}{app:matched-figure_05}
\resultplot{comparison}{figure_06}{SGD and Adam comparison (paper Figure~6): authors' and our clean, unaugmented CIFAR-10 CNN results with SGD and Adam. Authors' endpoints use full released histories; ours use 50,000 SGD updates or 400 Adam epochs.}{app:matched-figure_06}
\resultplot{comparison}{figure_07}{CIFAR-100 CNN width sweeps (paper Figure~7): authors' and our clean CIFAR-100 CNN results at eleven common widths. The authors' five-trial mean and standard deviation use full histories; ours use one seed and 50,000 SGD updates.}{app:matched-figure_07}
\resultplot{comparison}{figure_09}{Training trajectories and heatmaps (paper Figure~9): authors' and our 20\%-noise CIFAR-10 ResNet18 trajectories and width-by-time heatmaps. Ours cover 400 epochs; the author reconstruction retains its labeled released coverage.}{app:matched-figure_09}
\resultplot{comparison}{figure_10}{Training trajectories (paper Figure~10(c)): authors' and our width-128 CIFAR-10 CNN trajectories at the common 0 and 20\% noise settings. Authors' clean full-data and 50,000-example/20\%-noise references use saved measurement numbers; ours use recorded optimizer steps through 50,000 updates. Logging intervals and horizons are not matched. No smoothing is applied. The unresolved 10\%-noise author source is excluded.}{app:matched-figure_10}
\resultplot{comparison}{figure_11a}{Effect of training-set size (paper Figure~11(a)): authors' and our 20\%-noise CIFAR-10 CNN sample-size curves, using 12,500/25,000/50,000 examples and eleven common widths. Authors' endpoints use full recorded histories; ours stop at 50,000 updates. Our 10\%-noise subset runs are excluded because no matching author source was identified.}{app:matched-figure_11a}
\resultplot{comparison}{figure_12}{Sample-size heatmaps and slices (paper Figure~12): authors' and our 20\%-noise CIFAR-10 CNN sample-size heatmaps and slices. Both use three sample sizes, eleven common widths, and a common color scale; training horizons differ.}{app:matched-figure_12}
